一块盘坏了。在存着几 PB 数据的集群里，这算不上事故，只是日常——磁盘、整机乃至整个机架都会以可预测的频率退出服务。早年的对策很直接：每份数据存三副本，坏掉一份就从另外两份拷回来，存储成本则要乘以三。

现在 HDFS 与 Ceph 的纠删码配置换了一套算法。把一份数据切成 6 块，再算出 3 块校验块，9 块分散到 9 台机器上。任意 3 台同时失效，剩下的 6 块仍能把原数据完整还原。存储开销从 \(200\%\) 降到 \(50\%\)，容错能力反而提高了。多出来的 3 块不是任何一块原始数据的复制品，它们与那 6 块被同一个代数关系绑在一起，这个关系就是 Reed--Solomon 码。

所以真正的问题不是``要不要冗余''，而是冗余该怎么组织。多抄几遍只制造了副本，既没有说明副本之间应当满足什么，也没有告诉接收者版本冲突时该信谁。

本章的统一答案是几何的。把 \(k\) 个符号的消息映射成 \(n\) 个符号的码字，等于把 \(q^k\) 个消息稀疏地撒进一个有 \(q^n\) 个点的空间。消息原本挨得很紧，编码之后被人为拉开；噪声会把发出去的码字推离原位，只要推得不够远，接收者仍能看出它原本坐在哪里。

\begin{center}
\begin{tikzpicture}[x=1cm,y=1cm,font=\small,>=stealth]
  \draw[black!45,fill=blue!8] (0,0) rectangle (2.4,2.4);
  \foreach \i in {0,1,2,3,4,5}{
    \foreach \j in {0,1,2,3,4,5}{
      \fill[blue!70] (0.2+0.4*\i,0.2+0.4*\j) circle (1.3pt);
    }
  }
  \node[below,font=\footnotesize] at (1.2,-0.25) {\(q^{k}\) 个消息};
  \draw[->,thick] (2.9,1.2) -- (4.1,1.2);
  \node[above,font=\footnotesize] at (3.5,1.28) {编码};
  \draw[black!45,fill=blue!4] (4.6,-0.7) rectangle (10.8,3.1);
  \foreach \p/\q in {5.5/0.1, 7.1/1.7, 8.7/0.2, 9.9/2.1, 6.4/2.4, 9.5/0.7}{
    \draw[black!35,dashed] (\p,\q) circle (0.66);
    \fill[blue!70] (\p,\q) circle (1.9pt);
  }
  \draw[<->,black!70] (5.5,0.1) -- (7.1,1.7);
  \node[black!70,font=\footnotesize] at (6.65,0.62) {\(d\)};
  \node[below,font=\footnotesize] at (7.7,-0.95) {\(q^{n}\) 个串，其中只有 \(q^{k}\) 个是码字};
\end{tikzpicture}

\emph{编码不增加信息，只重新安排位置：同样多的消息被搬进一个大得多的空间，彼此之间空出了距离 \(d\)。虚线圈是每个码字的势力范围，噪声只要没把接收串推出圈外，归属就没有歧义。整章的判断几乎都能从这幅图里读出来。}
\end{center}

接下来几节各自补上这幅图的一块。有限域提供一个能做除法的离散数系；没有它，``解一个方程''这句话都站不住。多项式给出把距离拉开的最省办法：\(k\) 个点唯一确定一个 \(k-1\) 次多项式，多发几个点就自带冗余，而任意 \(k\) 个点都足以还原——Reed--Solomon 码与 Shamir 秘密共享共用的正是这一条。Hamming 距离把``能纠几个错''翻译成``那些圈会不会相交''。稀疏图码换了一种代价结构：用大量局部约束替换一个全局约束，让出一点可解码性，换回可以承受的计算量。

按依赖顺序排列，本章五篇分别是：

\begin{itemize}
\tightlist
\item \hyperref[sec:04-Discrete-Mathematics/05-Coding-and-Polynomial-Methods/01]{``有限域与多项式插值''}：为什么编码必须住在域上，以及低次多项式的刚性从何而来；
\item \hyperref[sec:04-Discrete-Mathematics/05-Coding-and-Polynomial-Methods/02]{``码距决定检测与纠错能力''}：把检错、纠错、纠擦除三种能力统一读成球是否相交；
\item \hyperref[sec:04-Discrete-Mathematics/05-Coding-and-Polynomial-Methods/03]{``Reed--Solomon 编码''}：多项式怎样把上一条几何判断变成可部署的方案；
\item \hyperref[sec:04-Discrete-Mathematics/05-Coding-and-Polynomial-Methods/04]{``稀疏图码与消息传递''}：全局约束换成局部约束之后，解码变成图上的迭代推断；
\item \hyperref[sec:04-Discrete-Mathematics/05-Coding-and-Polynomial-Methods/05]{``Fountain 码与无速率传输''}：当困难从``符号被改坏''变成``不知道丢了哪些包''时，编码思路怎样跟着变。
\end{itemize}

若只想解决眼前的一个问题，也不必顺读。想弄清冗余为什么真能纠错，第二篇够了；想要一个确定的、可以手工验证的方案，读第一、三篇；关心长码在现实成本下怎样解码，直接进第四篇，它用到的 Tanner 图建立在\hyperref[ch:028]{``图论：局部邻接怎样形成全局网络''}的二分图语言上。

读的时候始终盯住两个区分：信道造成的是错误、擦除还是软噪声？手上这条结论是存在性边界、确定性保证，还是高概率表现？把这两件事说清楚，比记住码名有用得多。至于可靠传输的极限究竟在哪里——不是某个具体的码能做到什么，而是任何码都越不过什么——那是信息论部分\hyperref[ch:058]{``有噪信道编码：噪声中为什么仍能可靠通信''}要回答的问题。
